With one validation candidate per grid cell, adding repair to both starts changes ordinary fidelity by 0.68 [-0.40, 1.95] percentage points for clean teachers and 0.40 [-0.19, 1.13] under noise (descriptive 95\% intervals in brackets). Both intervals include zero. Adding the pure start with repair gives 0.13 [-0.98, 1.07] and -0.32 [-1.40, 0.13] percentage points, respectively. Thus the positive mean repair effects are not stable enough in this follow-up to establish a clear advantage after controlling validation opportunities; the pure-start effect remains inconclusive. The original validation policy versus the one-per-cell full control differs by -0.16 [-0.46, 0.01] and 0.03 [-0.08, 0.13] percentage points. These controls leave training search and computational effort unequal. They do not redefine the deployed method or establish equivalence.

At the prespecified secondary sixteen-condition ceiling, the same repair contrast is 1.19 [0.08, 2.60] percentage points for clean teachers and 0.48 [-0.42, 1.44] under noise. The clean interval excludes zero, giving narrower evidence for repair at this larger ceiling; the noisy and pure-start effects remain inconclusive. This secondary result does not replace the eight-condition endpoint.
